% Part: normal-modal-logic
% Chapter: axioms-systems
% Section: provability-properties

\documentclass[../../../include/open-logic-section]{subfiles}

\begin{document}

\olfileid{nml}{prf}{prp}

\olsection{\usetoken{S}{derivability} चे गुणधर्म}

\begin{prop}\ollabel{prop:derivabilityfacts}
  $\Sigma$ ही मोडल प्रणाली आणि $\Gamma$ हा मोडल
  !!{formula}s चा संच असू द्या. पुढील गुणधर्म लागू होतात:
  \begin{enumerate}
  \item\ollabel{prop:derivabilityfacts-monotonicity} \emph{एकस्वनिकता}:
    $\Gamma \Proves[\Sigma] !A$ आणि
    $\Gamma \subseteq \Delta$ असेल, तर
    $\Delta \Proves[\Sigma] !A$;
  \item\ollabel{prop:derivabilityfacts-reflexivity}
    \emph{परावर्तिता}: $!A \in \Gamma$ असेल, तर $\Gamma
    \Proves[\Sigma] !A$;
  \item\ollabel{prop:derivabilityfacts-cut} \emph{कट}: $\Gamma
    \Proves[\Sigma] !A$ आणि $\Delta \cup \{!A\} \Proves[\Sigma] !B$
    असेल, तर $\Gamma \cup \Delta \Proves[\Sigma] !B$;
  \item\ollabel{prop:derivabilityfacts-deduction} \emph{निगमन प्रमेय}: $\Gamma \cup \{!B\}
    \Proves[\Sigma] !A$ असेल तेव्हा आणि तेव्हाच $\Gamma \Proves[\Sigma] !B \lif
      !A$;
  \item \ollabel{prop:derivabilityfacts-ruleT} \emph{नियम T}: जर
    $\Gamma \Proves[\Sigma] !A_1$ आणि \dots आणि $\Gamma
    \Proves[\Sigma] !A_n$ असेल आणि $!A_1 \to (!A_2 \lif \cdots (!A_n \lif
    !B)\cdots)$ हा सर्वतः-सत्य दृष्टांत असेल, तर $\Gamma
    \Proves[\Sigma] !B$.
  \end{enumerate}
\end{prop}

सिद्धता हा सोपा सराव आहे.
\olref{prop:derivabilityfacts} मधील
\olref{prop:derivabilityfacts-ruleT} या भागावरून,
उदाहरणार्थ, $\Gamma \Proves[\Sigma] !A \lor !B$
आणि $\Gamma \Proves[\Sigma] \lnot !A$ असेल, तर
$\Gamma \Proves[\Sigma] !B$ मिळते.
तसेच पुढे $\Gamma \cup \{ !A\} \Proves[\Sigma] !B$
ऐवजी $\Gamma, !A \Proves[\Sigma] !B$ असे लिहू.

\begin{defn}
  $\Gamma \Proves[\Sigma] !A$ असेल तेव्हा नेहमी
  $!A \in \Gamma$ असेल, तर आणि तेव्हाच संच
  $\Gamma$ हा प्रणाली~$\Sigma$ च्या सापेक्ष
  \emph{निगामीदृष्ट्या बंद} आहे, असे म्हणतो.
\end{defn}

\end{document}
